% ==========================================================
% 方格最短路径与相遇概率
% 难度：★★★ 难题
% ==========================================================

\newcommand{\qProbPermGridPathMeetStem}{%
在一个 $3\times 3$ 的方格道路网中，$M$ 为左下角，$N$ 为右上角。甲从 $M$ 出发前往 $N$，乙从 $N$ 出发前往 $M$。两人各自等可能地随机选择一条最短路径，以相同速度同时出发。反对角线上有四个交汇点 $A_1,A_2,A_3,A_4$。判断下列结论：\\
A：甲从 $M$ 到 $N$ 有 $120$ 种方法\\
B：甲必须经过 $A_2$ 到达 $N$ 有 $64$ 种方法\\
C：甲、乙在 $A_2$ 相遇的概率为 $\frac{81}{400}$\\
D：甲、乙相遇的概率为 $\frac{1}{2}$
}

\newcommand{\qProbPermGridPathMeetAnswer}{%
C
}

\newcommand{\qProbPermGridPathMeetSolution}{%
\textbf{判断 A：总最短路径数}\\
甲从 $M$ 到 $N$，共走六步：三步向右，三步向上。\\
只需从六步中选择三步向右：$C_6^3=20$。A 错。

\textbf{判断 B：经过 $A_2$ 的路径数}\\
从 $M$ 到 $A_2$ 需要：一步向右，两步向上。\\
方法数为 $C_3^1=3$。\\
从 $A_2$ 到 $N$ 需要：两步向右，一步向上。\\
方法数为 $C_3^1=3$。\\
经过 $A_2$ 的最短路径数为 $3\cdot 3=9$。B 错。

\textbf{判断 C：在 $A_2$ 相遇}\\
甲经过 $A_2$ 有 $9$ 条路径。\\
由对称性，乙经过 $A_2$ 也有 $9$ 条路径。\\
因 $A_2$ 距离双方起点都恰好为三步，两人同速同时出发，\\
只要路径都经过 $A_2$，就会在第三步结束时相遇。\\
有利路径对数为 $9^2=81$。总路径对数为 $20^2=400$。\\
所以 $P(\text{在 }A_2\text{ 相遇}) = \frac{81}{400}$。C 正确。

\textbf{判断 D：两人相遇的总概率}\\
两人相遇只能在距离起点均为三步的反对角线上。
\begin{itemize}
    \item 在 $A_1$ 相遇：双方路径唯一，共 $1$ 个路径对。
    \item 在 $A_2$ 相遇：共 $81$ 个路径对。
    \item 在 $A_3$ 相遇：由对称性，共 $81$ 个路径对。
    \item 在 $A_4$ 相遇：双方路径唯一，共 $1$ 个路径对。
\end{itemize}
相遇概率为 $\frac{1+81+81+1}{400} = \frac{41}{100}$。D 错误。

综上，正确选项为 C。
}

\newcommand{\qProbPermGridPathMeetTeachingNotes}{%
\textbf{【避坑与边界说明】}
\begin{itemize}
    \item 选定三步向右的位置后，另外三步自动向上，\\
    不应再重复计数。
    \item “经过某点”应拆成起点到该点、该点到终点\\
    两个阶段，利用分步乘法。
    \item 经过某点不一定相遇，必须步数相同。\\
    本题因同速同时出发，相遇只能在第三步结束。
    \item 双方路径选择独立，总样本空间为 $20^2 = 400$。
\end{itemize}
}
